% IEEE T-ITS manuscript — experimental results of the TriAction study.
% Every number in this file is traceable to
% doc/experiment_report_no_early_stopping_2026-07-24.md and the saved
% predictions under logs/train/K_*/; nothing is estimated or rounded beyond
% the reported precision.
\documentclass[journal]{IEEEtran}

\usepackage{amsmath,amssymb}
\usepackage{booktabs}
\usepackage{multirow}
\usepackage{graphicx}
\usepackage[colorlinks=false,hidelinks]{hyperref}

\newcommand{\ie}{i.e.}
\newcommand{\eg}{e.g.}
\newcommand{\micro}{micro-accuracy}
\newcommand{\mF}{macro-F$_1$}

\begin{document}

\title{Multi-View Recognition of Driver Head Movements:\\
A Measurement-First Study with Person-Held-Out Evaluation}

\author{Kaixu~Chen,~\IEEEmembership{}%
\thanks{K.~Chen is with the University of Tsukuba, Tsukuba, Japan
(e-mail: chenkaixusan@gmail.com).}%
\thanks{Manuscript draft generated from the 2026-07 experimental campaign;
co-author list and funding to be completed.}}

\markboth{IEEE Transactions on Intelligent Transportation Systems}%
{Chen: Multi-View Recognition of Driver Head Movements}

\maketitle

\begin{abstract}
We study four-class recognition of driver head-movement direction
(left/right/down/up) from three synchronized cabin cameras, and report the
results of a systematic 100+-run campaign on 22 drivers.  Our contributions
are threefold.  (1)~\emph{Method}: a parameter-efficient multi-view
architecture that keeps a single pretrained video transformer frozen and
learns only a small cross-view token-attention module (11M trainable
parameters).  Under person-held-out evaluation it improves \mF{} by
$+6.3$ points over an equally pretrained single-view model
($p{=}1.9\times10^{-4}$, paired bootstrap), and under a nested,
selection-free protocol its \micro{} advantage grows to $+5.1$ points
($p{=}8.6\times10^{-12}$), while halving the across-driver variance.
(2)~\emph{Measurement}: we show that two common protocol defects---a video-level
split that leaks driver identity, and per-batch averaging of epoch
metrics---jointly fabricated a coherent but false architecture narrative on
this task; several ``validated'' conclusions, including our own earlier ones,
reversed under a sound protocol.  (3)~\emph{Negative results}: an analytic
head-pose feature stream carries genuine signal in isolation, yet no injection
route (token-, query-, or logit-level) survived robustness checks, and five
class-imbalance remedies failed to prevent rare-class collapse.  We release
the person-held-out protocol, statistical-testing tooling, and all per-run
predictions to support measurement-first research on small-cohort driver
monitoring.
\end{abstract}

\begin{IEEEkeywords}
Driver monitoring, head movement recognition, multi-view fusion, video
transformers, evaluation protocols, negative results.
\end{IEEEkeywords}

\section{Introduction}
\IEEEPARstart{D}{river}-monitoring systems infer attention from head and gaze
behaviour.  We consider a cabin setup with three synchronized cameras (front,
left, right) and the task of classifying short clips into four head-movement
directions.  The setting is representative of practical driver-monitoring
research: a small cohort (22 drivers), heavy class imbalance, and strong
inter-driver appearance variation.

This paper is deliberately \emph{measurement-first}.  Our initial campaign,
run under a conventional protocol (random video-level split, standard metric
logging), produced a seemingly well-supported progression of architectural
improvements.  Auditing the protocol invalidated essentially all of it.  We
therefore report (i)~the method that survives a person-held-out, statistically
tested evaluation; (ii)~a controlled account of how the defective protocol
manufactured false conclusions; and (iii)~the negative results of our attempts
to exploit head-pose priors and to fix rare-class collapse.

\section{Task and Data}
\label{sec:data}
\subsection{Recording and annotation}
Each of 22 drivers was recorded in four illumination conditions
(day/night $\times$ high/low), yielding 88 sessions.  Three synchronized
cameras (front, left, right; $332{\times}224$\,px) observe the cabin.
Timeline annotations mark head-movement segments with one of eight direction
labels, merged to four (left/right/down/up; vertical components take
precedence).  Three annotators labelled every session; frame-level
inter-annotator agreement is 91.0\% overall and 94\% for the
\emph{right} class, so label noise does not explain the failure modes reported
below.  Segment-level class frequencies are imbalanced: down 43\%, left 28\%,
right 21\%, up 9\% (pooled over the five held-out folds).  SAM-3D body
keypoints (70 points, 2-D and 3-D) are available for every view.

\subsection{Evaluation protocol}
\label{sec:protocol}
\textbf{Person-held-out 5-fold CV.}  Drivers---not videos---are the unit of
splitting: the 22 drivers are partitioned into five folds; each fold's
sessions form the evaluation set of a model trained on the remaining drivers.
Pooling the folds scores every driver exactly once with a model that never saw
them ($n{=}3954$ segments).  \textbf{Nested variant.}  For selection-free
reporting, three additional drivers are carved out of each training set as an
inner validation set; checkpoints are selected on the inner set and reported
on the held-out fold, so the reported data never influences model selection.
\textbf{Metrics.}  \micro{} (segment accuracy), macro-accuracy, and \mF{},
accumulated over the full epoch (Sec.~\ref{sec:pitfalls}).  \textbf{Statistics.}
All model comparisons are paired on identical segment sets (verified
fold-by-fold): McNemar's test with continuity correction and a paired
bootstrap (1{,}500--2{,}000 resamples, 95\% CIs).  The pooled binomial 95\% CI
is $\pm1.6$ points; a single 8-video split, by contrast, carries $\pm6.4$
points and up to 15.6 points of across-fold spread, which is why no
single-split number appears in this paper.

\subsection{Implementation}
All arms train 50 epochs without early stopping (Adam, lr $10^{-4}$, cosine
schedule, effective batch 16, bf16), 32 uniformly sampled frames at
$224^2$ for ViViT arms, on one H100 per run.  The video backbone is
ViViT-B/16x2 pretrained on Kinetics-400.  Class-frequency weights are applied
to the cross-entropy loss in all arms.

\section{Method}
\label{sec:method}
The method that survives evaluation is deliberately minimal
(Fig.~\ref{fig:arch}; ``MV-Min'' below).  A \emph{single, frozen} pretrained
ViViT encoder processes each view.  Per view, the $3137$ output tokens are
compressed to $17$: the CLS token plus one spatially averaged token per
temporal tubelet.  A learned additive view embedding tags each view's tokens,
the three token sets are concatenated, and a 2-layer, 8-head transformer whose
attention spans all views performs cross-view fusion.  Mean pooling and a
linear head produce the logits.  Only the fusion module and head train
(11M parameters, 4\% of one finetuned backbone; the backbone runs under
\texttt{no\_grad}).

Design choices that our ablations showed to be unnecessary---partial backbone
unfreezing, keypoint-trajectory tokens, keypoint-query pooling, gated
aggregation---are excluded from the method and revisited as ablations in
Sec.~\ref{sec:ablation} and negative results in Sec.~\ref{sec:keypoints}.

\begin{figure}[t]
\centering
\framebox[\linewidth]{\parbox{0.95\linewidth}{\centering\vspace{18mm}
\textit{[Figure: three views $\rightarrow$ shared frozen ViViT $\rightarrow$
17 tokens/view + view embeddings $\rightarrow$ cross-view transformer
$\rightarrow$ mean pool $\rightarrow$ head]}\vspace{18mm}}}
\caption{MV-Min: a frozen shared ViViT encoder with a small trainable
cross-view token-attention module.}
\label{fig:arch}
\end{figure}

\section{Main Results}
\label{sec:main}

\begin{table}[t]
\caption{Person-held-out 5-fold CV (pooled $n{=}3954$; seed 42).
``Spread'' is the across-fold standard deviation of \micro.}
\label{tab:main}
\centering
\begin{tabular}{lcccc}
\toprule
Model & \micro & mac.\ acc & \mF & spread \\
\midrule
Majority class            & 42.7 & 25.0 & 15.0 & --   \\
3D-CNN, single view       & 47.0 & 31.1 & 26.8 & .156 \\
TS-CVA mid fusion (3D-CNN)& 45.4 & 31.9 & 28.1 & --   \\
\midrule
ViViT, single view        & 52.5 & 35.5 & 29.8 & .134 \\
ViViT, 3-view late fusion & 52.6 & 35.7 & 31.9 & .114 \\
\textbf{MV-Min (ours)}    & 51.5 & 38.0 & \textbf{36.5} & \textbf{.090} \\
\bottomrule
\end{tabular}
\end{table}

\begin{table}[t]
\caption{Paired comparisons on identical segments ($n{=}3954$).
$\Delta$ in points with 95\% bootstrap CI; McNemar $p$ for \micro.}
\label{tab:paired}
\centering
\begin{tabular}{lccc}
\toprule
Comparison & $\Delta$\micro & $\Delta$\mF & $p$ (McNemar) \\
\midrule
MV full vs.\ 3D-CNN single & $+5.1\,[+3.6,+6.6]$ & $+9.2$ & $1.9{\times}10^{-11}$ \\
MV full vs.\ ViViT single  & $-0.4\,[-1.7,+0.9]$ & $+6.3\,[+4.8,+7.7]$ & $0.55$ \\
MV full vs.\ ViViT late    & $-0.4\,[-1.8,+0.9]$ & $+4.1\,[+2.4,+5.8]$ & $0.55$ \\
\bottomrule
\end{tabular}
\end{table}

Table~\ref{tab:main} and the paired tests in Table~\ref{tab:paired} support
three statements.  \textbf{(1) Pretraining dominates \micro.}  Replacing the
3D-CNN with a pretrained ViViT adds $+5.5$ \micro{} points; no fusion
mechanism adds more on top of it.  \textbf{(2) Cross-view attention buys
balance and robustness, not raw accuracy.}  Against the equally pretrained
single-view and trivial late-fusion controls, our fusion is \micro-neutral but
gains $+6.3$ and $+4.1$ \mF{} respectively (both significant) and cuts the
across-driver spread from .134 to .090.  On a task whose practical failures
are rare-class collapse and driver variance, this is the relevant axis.
\textbf{(3) A previously ``winning'' fusion reverses.}  TS-CVA mid fusion,
the best fusion under our earlier leaky protocol (48.9 vs.\ 37.2 single-view),
scores \emph{below} the single-view 3D-CNN under person hold-out
(45.4 vs.\ 47.0).

\subsection{Nested, selection-free evaluation}
Selecting checkpoints on the same people one reports on inflates results
(measured: $+2.0$ points for the 3D-CNN arm, $+0.9$ for the multi-view arm).
Under the nested protocol of Sec.~\ref{sec:protocol} every number drops, but
the ordering sharpens in our favour: MV-Min is the best arm (50.6 \micro),
and the paired advantage of cross-view attention over the single-view control
\emph{grows} to $+5.1$ points ($p{=}8.6\times10^{-12}$); \mF{} $+2.9$
($[+1.4,+4.4]$).  The stricter the protocol, the clearer the method's benefit
---the opposite of what an overfit result would show.

\subsection{Seed variance}
Replicating both headline arms with seeds 42/43/44: single-view ViViT
52.5/53.0/51.9 \micro; the full multi-view arm 52.1/52.6/51.9.  Run-to-run
variation ($\pm0.6$) is well inside the protocol CI, and the \mF{} advantage
of cross-view fusion persists in all three seeds (mean $+5.0$).

\section{Attribution Ablation}
\label{sec:ablation}

\begin{table}[t]
\caption{Component attribution under person-held-out CV
(paired vs.\ the row above it or as indicated).}
\label{tab:ablation}
\centering
\begin{tabular}{lccc}
\toprule
Arm & \micro & \mF & paired effect \\
\midrule
MV-Min (frozen, no kpt)      & 51.5 & 36.5 & reference \\
\;+ partial unfreeze (4)     & 52.4 & 36.2 & n.s.\ / n.s. \\
\;+ kpt-query pooling        & 52.1 & 36.0 & $+1.3^{*}$ / n.s. \\
\;+ unfreeze depth 6         & 50.5 & 35.5 & n.s. \\
\;+ mirror augmentation      & 53.0 & 35.3 & n.s. \\
\bottomrule
\multicolumn{4}{l}{\footnotesize $^{*}$McNemar $p{=}0.004$; not robust under
the nested protocol.}
\end{tabular}
\end{table}

Table~\ref{tab:ablation}: no addition to the minimal form produces a robust
paired improvement on both metrics.  Notably, conclusions that appeared firmly
established under the leaky protocol---``query pooling beats token injection
by $+3.5$'', ``unfreeze depth 6 is optimal'', ``partial unfreezing gains
$+2.9$''---all shrink to non-significance or reverse.  This motivates the
minimal form as the published method.

\section{Head-Pose Priors: Signal Without a Robust Injection Route}
\label{sec:keypoints}
Head movements should be readable from head pose, and our scenario provides
SAM-3D keypoints.  A diagnosis of why naive keypoint streams failed is
instructive: the raw $70{\times}3$ tensor is 60\% hand rows (moving $4\times$
more than the head), in unnormalized coordinates that encode driver identity,
compressed into one token competing with 51 video tokens.  The head signal
never reached the model.

We therefore computed analytic features: 8 head-related points, a per-frame
head-local frame from the ear axis and nose direction, yaw/pitch/roll plus
first differences and centroid velocity plus local-frame shape (33-D/frame),
with per-segment mean-pose subtraction to remove resting posture (identity).

\textbf{The signal is real.}  A 64-unit MLP on these features alone---no
video---reaches 27.2 \mF{} on unseen drivers with \emph{non-zero F$_1$ on
every class}: right 0.24 and up 0.12, where every video model scores
${\approx}0.03$ and $0.00$ (Table~\ref{tab:collapse}).

\textbf{No injection route survives.}  Per-frame pose tokens initially gave
the best \micro{} of the campaign (53.3, $+1.7$ over MV-Min,
$p{=}8.9\times10^{-5}$), but three robustness checks demoted it:
(i)~seed sensitivity ($+1.7/+3.3/+0.4$; significant in 2 of 3 seeds);
(ii)~\emph{reversal} under the nested protocol ($-2.5$,
$p{=}3.2\times10^{-7}$)---part of the gain came from checkpoint selection on
the scored drivers; (iii)~a representation control showed raw flattened
keypoint tokens capture most of the effect ($+1.1$, $p{=}0.021$), leaving
$+0.6$ (n.s.)\ for the analytic encoding.  Logit-level fusion of the pose
probe with the video model, pre-registered as a 50/50 average, degraded
\micro{} by $9.1$ points ($p{=}5.4\times10^{-30}$): the probe is 22 points
weaker overall, and the equal-weight blend loses more on frequent classes
than it recovers on rare ones.  We report these as negative results; the open
problem is an injection mechanism that preserves the probe's rare-class
competence.

\section{Rare-Class Collapse and Failed Remedies}
\label{sec:collapse}

\begin{table}[t]
\caption{Per-class F$_1$ under person-held-out CV: every video model
collapses onto the frequent classes; only the pose-only probe predicts all
four.}
\label{tab:collapse}
\centering
\begin{tabular}{lcccc}
\toprule
Model & left & right & down & up \\
\midrule
MV-Min                  & .538 & .173 & .673 & .078 \\
\;+ pose tokens         & .532 & .179 & .695 & .015 \\
Pose-only probe (33-D)  & .327 & .237 & .409 & .117 \\
\bottomrule
\end{tabular}
\end{table}

\begin{table}[t]
\caption{Remedies for class collapse, all evaluated under the sound
protocol.  None succeeded.}
\label{tab:remedies}
\centering
\begin{tabular}{lcc}
\toprule
Remedy & \micro & outcome \\
\midrule
Class-weighted CE (default)      & 51.5 & collapse persists \\
Focal loss                       & $-5.4$ vs.\ ref & harmful \\
Balanced softmax                 & 26.4 & catastrophic \\
Clip-mean subtraction            & 43.6 & harmful \\
Pose-token fusion                & 53.3$^{\dagger}$ & up-F$_1$ drops \\
Probe--model logit fusion (50/50)& 44.1 & harmful \\
\bottomrule
\multicolumn{3}{l}{\footnotesize $^{\dagger}$not robust (Sec.~\ref{sec:keypoints}).}
\end{tabular}
\end{table}

Two persistent failure modes define the real difficulty of this task.
\textbf{Collapse}: \emph{right} and \emph{up} are essentially never predicted
by any video model (Table~\ref{tab:collapse}), despite 94\% annotator
agreement on \emph{right}.  Six remedies failed (Table~\ref{tab:remedies}).
\textbf{Driver memorisation}: under person hold-out the best checkpoint lands
at epoch 0--4 in every arm---models start memorising training drivers almost
immediately, consistent with the large across-driver spread.  Removing static
appearance (clip-mean subtraction) removed useful signal as well.  We consider
identity-invariant representations the key open problem for small-cohort
driver monitoring.

\section{How a Defective Protocol Fabricated a Narrative}
\label{sec:pitfalls}
Our first 40+ runs used (a)~a random \emph{video-level} split---all seven
validation drivers also appeared in training---and (b)~torchmetrics objects
evaluated per batch (size 2) and averaged, which at that batch size inflates
\mF{} by up to 17 points (0.52 logged vs.\ 0.35 true) and was inconsistently
applied across trainer families.  Together these produced a coherent, mutually
confirming, and false progression (54.0 $\rightarrow$ 56.9 $\rightarrow$
58.3 ``SOTA'' with supporting ablations).  Under the sound protocol:
the ``$+3.5$ query-pooling'' effect vanished; ``optimal unfreeze depth''
became n.s.; the TS-CVA fusion champion fell below its own single-view
baseline; and an early ``architectures are indistinguishable'' reading---drawn
from the same leaky split---also proved wrong.  Single-split evaluation on
8 videos carries a $\pm6.4$-point CI and 15.6 points of fold spread; every
effect we chased for four rounds was smaller than that noise.  We include
this account because both defects are common, silent, and jointly
self-reinforcing; the protocol of Sec.~\ref{sec:protocol} is released as the
antidote.

\section{Conclusion}
Under person-held-out, statistically tested evaluation of 20+ arms
(100+ trained models), the defensible findings are:
(1)~Kinetics-pretrained video transformers dominate \micro{} on this task
($+5.5$ over a 3D-CNN);
(2)~a frozen-backbone cross-view token-attention module---4\% of the
trainable parameters of one finetuned backbone---adds class balance
($+6.3$ \mF), across-driver robustness (spread halved), and, under the
strictest selection-free protocol, $+5.1$ \micro{}
($p{=}8.6\times10^{-12}$) over its single-view control;
(3)~head-pose priors carry real, complementary signal but lack a robust
injection route;
(4)~rare-class collapse and driver memorisation, not architecture, are the
binding constraints---six remedies failed;
(5)~two common measurement defects sufficed to fabricate an entire false
research narrative on a small cohort.
We release the protocol, tests, and per-run predictions.

\section*{Reproducibility}
All splits are person-disjoint and seeded; every comparison is paired on
identical segments with McNemar and bootstrap statistics; per-run predictions
are archived.  Code, job scripts (90 runs), and the audit trail accompany the
paper.

\begin{thebibliography}{1}
\bibitem{vivit} A.~Arnab \emph{et al.}, ``ViViT: A video vision transformer,''
in \emph{Proc. ICCV}, 2021.
\bibitem{kinetics} W.~Kay \emph{et al.}, ``The Kinetics human action video
dataset,'' \emph{arXiv:1705.06950}, 2017.
\bibitem{mcnemar} Q.~McNemar, ``Note on the sampling error of the difference
between correlated proportions or percentages,'' \emph{Psychometrika}, 1947.
% TODO: related work on driver monitoring / multi-view action recognition.
\end{thebibliography}

\end{document}
